% Front cover for "Một số thuật ngữ trong Toán rời rạc". Compile with xelatex.
\documentclass[11pt]{article}
\usepackage[paperwidth=8.5in,paperheight=11in,margin=0pt]{geometry}
\usepackage{fontspec}
\setmainfont{TeX Gyre Termes}
\usepackage{tikz}
\usepackage{xcolor}
\definecolor{ink}{HTML}{0F2A43}
\definecolor{accent}{HTML}{2C7A7B}
\definecolor{gold}{HTML}{C9A227}
\definecolor{paper}{HTML}{F6F8FA}
\pagestyle{empty}
\setlength{\parindent}{0pt}
\begin{document}
\begin{tikzpicture}[x=1in,y=1in]
  % background
  \fill[paper] (0,0) rectangle (8.5,11);
  % top ink band + accent rule
  \fill[ink] (0,7.5) rectangle (8.5,11);
  \fill[accent] (0,7.36) rectangle (8.5,7.5);
  \fill[gold]   (0,7.30) rectangle (8.5,7.34);
  % title
  \node[anchor=west, text=white, align=left] at (0.85,9.85)
    {\fontsize{31}{36}\selectfont Một số thuật ngữ};
  \node[anchor=west, text=white, align=left] at (0.85,9.15)
    {\fontsize{31}{36}\selectfont trong Toán rời rạc};
  \node[anchor=west, text=gold] at (0.85,8.45)
    {\fontsize{16}{19}\selectfont\itshape Từ điển chú giải Anh\,–\,Việt};
  \node[anchor=west, text=paper] at (0.85,8.02)
    {\fontsize{11}{13}\selectfont Selected Terms in Discrete Mathematics};
  % Petersen-graph motif
  \begin{scope}[shift={(4.25,4.45)}]
    \foreach \i in {0,...,4} {
      \coordinate (O\i) at ({90+72*\i}:1.75);
      \coordinate (I\i) at ({90+72*\i}:0.9);
    }
    \draw[accent, line width=1.1pt] (O0)--(O1)--(O2)--(O3)--(O4)--cycle;
    \foreach \i in {0,...,4} \draw[accent!75, line width=0.9pt] (O\i)--(I\i);
    \draw[accent!75, line width=0.9pt] (I0)--(I2)--(I4)--(I1)--(I3)--cycle;
    \foreach \i in {0,...,4} {
      \fill[ink] (O\i) circle (3.2pt);
      \fill[gold] (I\i) circle (3.0pt);
    }
  \end{scope}
  % bottom band
  \fill[accent] (0,1.5) rectangle (8.5,1.56);
  \node[anchor=west, text=ink] at (0.85,1.12)
    {\fontsize{14}{16}\selectfont Duc A. Hoang};
  \node[anchor=west, text=ink!72] at (0.85,0.74)
    {\fontsize{10}{12}\selectfont VNU University of Science \;\textbullet\; 2026 \;\textbullet\; CC BY-SA 4.0};
\end{tikzpicture}
\end{document}
